% 中英两份文档共用的排版头。语言相关的只有四处名称与正文中文字体，四处名称在下面
% 写成占位符，由 build/build.sh 逐份 sed 成本次的取值（占位符的字面写法见该脚本里
% 的 sed 那一行——这里不重复写，否则 sed 会连注释一起替换掉）：
%   页眉右端的手册名（PDF 标题也用它）
%   目录标题：「目录」/「Contents」
%   图注编号格式：「图#2. 」/「Figure #2. 」
%   正文里引用图时用的词：「图」/「Figure」
% 中文专有的断行规则与中文字体在 a3_cjk.tex 里，只对中文版追加。

\usepackage{caption}
\usepackage{etoolbox}
\usepackage{fancyhdr}
\usepackage{titlesec}
\usepackage{colortbl}
\usepackage{float}
\usepackage{tikz}
\usetikzlibrary{calc}

% 行距 ≈ 1.5×正文（正文 11pt → 基线约 16.5pt，与 A3 一致）
\linespread{1.2}
\setlength{\parskip}{5pt}
\setlength{\parindent}{0pt}
\setlength{\emergencystretch}{4em}
\raggedbottom
% 段落拆页防孤行：禁止一页顶部只余段落末行（widow）或一页底部只有段落首行（club），
% 避免出现“术支持。”孤行字等；CJK 拆行同样受此约束。
\widowpenalty=10000
\clubpenalty=10000
\displaywidowpenalty=10000

% ===== 标题：章 14 / 节 12 / 小节~10.9，思源宋体加粗（A3 实测值）=====
% 编号由 LaTeX 自动给出（目录才可“编号独立成列”对齐 A3），格式与参考一致。
\titleformat{\section}{\fontsize{14.3}{20}\selectfont\bfseries}{\thesection}{0.35em}{}
\titleformat{\subsection}{\fontsize{12}{16}\selectfont\bfseries}{\thesubsection}{0.35em}{}
\titleformat{\subsubsection}{\fontsize{10.9}{15}\selectfont\bfseries}{\thesubsubsection}{0.35em}{}
\titleformat{\paragraph}{\fontsize{10.9}{15}\selectfont\bfseries}{\theparagraph}{0.35em}{}
% 标题前后留白对齐 A3 参考（实测视觉留白：小节/小小节 上方≈28pt、下方≈18.6pt）
\titlespacing*{\section}{0pt}{8pt}{13pt}
\titlespacing*{\subsection}{0pt}{21pt}{14pt}
\titlespacing*{\subsubsection}{0pt}{21pt}{14pt}
\titlespacing*{\paragraph}{0pt}{10pt}{6pt}
% 正文标题编号需到第 4 级（如 1.3.3.1）；目录仅列到第 3 级（tocdepth=3，下方设置）
\setcounter{secnumdepth}{4}
% 标题防孤行：标题若落在页末某阈值内，先分页再排，避免标题成为页末孤行。
% 阈值按手册分流：开发手册默认 120pt（保持 28 页；附录章开篇标题除外）；
% 产品手册经实测提高到 160pt，防止“标题＋后随大图浮体”被顶成页末孤标题。
\makeatletter
\newlength{\aifc@orphanreserve}
\setlength{\aifc@orphanreserve}{120pt}
\newcommand{\aifc@needheadspace}{%
  \ifdim\pagetotal>\dimexpr\pagegoal-\aifc@orphanreserve\relax
    % \pagetotal 在“刚溢出、页面还没换出去”的瞬间是**旧页**的累计高度（会把已经
    % 判给下一页的那几行也算进来）。这时再 \newpage 就会把新页在几行之内截断，
    % 留下一整页空白。\pagetotal 不小于 \pagegoal 说明当前页本来就放不下了，
    % 分页由页面构造器自然处理，这里不要再插一脚。
    \ifdim\pagetotal<\pagegoal
      \newpage
    \fi
  \fi}
\pretocmd{\section}{\aifc@needheadspace}{}{\relax}
\pretocmd{\subsection}{\aifc@needheadspace}{}{\relax}
\pretocmd{\subsubsection}{\aifc@needheadspace}{}{\relax}
\makeatother
\setcounter{tocdepth}{3}
\renewcommand{\contentsname}{@CONTENTSNAME@}
% ===== 目录页 =====
% 章级目录条目列宽对齐 A3：编号左缘在 x43（缩进 0em），标题自 x55 起（编号盒 1.1em）；
% 二级条目沿用 article 类默认缩进 1.5em、编号盒 2.3em（实测编号 x59、标题 x84）。
\makeatletter
% 目录条目间距：章级与二级**用同一个 6pt**，让全篇条目等距。
% 房式原件只在章级给 11.3pt、二级不给，条目间距分两档（章前 ≈27pt）；
% 本教程章少、小节也少，分档留白会让目录看起来忽松忽紧，故两级统一。
% 实测：不给 \addvspace 的条目间距 = 15.58pt（纯基线距），加了 6pt → 21.56pt。
% 要还原 SE 观感，把下面两个 6pt 分别改回 11.3pt 与 0pt（删掉 \addvspace）。
% 编号列 0em 缩进 + 1.1em 编号盒 → 编号 x43、标题 x55。
\renewcommand*\l@section{\addvspace{6pt}\@dottedtocline{1}{0em}{1.1em}}
\renewcommand*\l@subsection{\addvspace{6pt}\@dottedtocline{2}{1.5em}{2.3em}}
\makeatother
% “目录”标题下方留白（首行再叠加大节分组留白后总距 ≈48.7pt）
\let\aifcstdTOC\tableofcontents
\renewcommand{\tableofcontents}{%
  \begingroup
    \titlespacing*{\section}{0pt}{10pt}{23pt}%
    \aifcstdTOC
  \endgroup}

% ===== 图注：小号居中 "图N. 说明" =====
% 中文版的「图1.」里那个空格是 xeCJK 自动加的（汉字与拉丁字母之间），
% 英文版没有这一层，所以英文的编号格式里要自己写一个空格。
\DeclareCaptionLabelFormat{a3}{@FIGLABEL@}
\DeclareCaptionFont{a3}{\fontsize{8.5}{11}\selectfont}
\renewcommand{\figurename}{@FIGNAME@}
\captionsetup[figure]{font={a3,rm},labelfont={a3,rm},labelformat=a3,labelsep=none,skip=3pt}

% ===== 表格：booktabs 灰色横线（A3：无竖线、灰线、上下粗中间细）=====
% 灰白斑马纹底色：正文第 2、4、6…行 \rowcolor{zebra}（zebra.py 注入）
\definecolor{zebra}{gray}{0.937}
\renewcommand{\heavyrulewidth}{1.6pt}
\renewcommand{\lightrulewidth}{0.55pt}
\arrayrulecolor[gray]{0.6}
\AtBeginEnvironment{longtable}{\small}
\AtBeginEnvironment{table}{\small}
\AtBeginEnvironment{verbatim}{\footnotesize}
\AtBeginEnvironment{Shaded}{\footnotesize}

% ===== 页眉页脚（仿 A3）=====
\fancypagestyle{frontmatter}{%
  \fancyhf{}%
  \fancyfoot[L]{NEXFORM ROBOTICS}%
  \fancyfoot[R]{\thepage}%
  \renewcommand{\headrulewidth}{0.4pt}%
  \renewcommand{\footrulewidth}{0.4pt}%
}
\fancyhf{}
\fancyhead[R]{@DOCSHORT@}
\fancyfoot[L]{NEXFORM ROBOTICS}
\fancyfoot[R]{\thepage}
\renewcommand{\headrulewidth}{0.4pt}
\renewcommand{\footrulewidth}{0.4pt}
\pagestyle{fancy}
\AtBeginDocument{\hypersetup{pdftitle={@DOCSHORT@}}}
% 目录/正文超链接改为黑色（去除默认红色链接色）
\AtBeginDocument{\hypersetup{colorlinks=true,linkcolor=black,urlcolor=black,citecolor=black,filecolor=black}}

% ===== 图片控制（居中、统一宽度上限）=====
% 版心实测 520.7pt × 651.6pt（a4paper, left=15mm right=12mm top=38mm bottom=30mm）。
% 全篇图统一按 0.70\linewidth = 364.5pt 宽，高上限 0.52\textheight = 338.8pt。
% 0.70 是实测出来的折中点：正文照片已统一裁成约 4:3 横构图（见
% build/compose_figures.py 的 CROPS），1.33:1 的图在此宽度下高约 273pt，
%   - 两图 + 两条图注 ≈ 574pt，还能剩 ~78pt 排正文 → 一页放得下两组，最省页数；
%   - 0.62\linewidth（322.8pt 宽／242pt 高）更省页但图明显偏小，读接口文字吃力；
%   - 0.82\linewidth（427pt 宽／320pt 高）图最大，但一页只放得下一张图，
%     页数从 15 页涨到 19 页、页底留白翻倍，反而不均匀。
% 高上限 338.8pt 对 4:3 图不生效（273 < 339），保留是为了防止日后换入更高的图时
% 撑破版心；更宽的整图（立柱、脚轮、遥操那三张 2.2~2.4:1）本就按宽度定尺寸。
% 代价：就地排版（见下）时整张图放不进当前页剩余空间就会整页下移、在页底留白。
\AtBeginDocument{\setkeys{Gin}{width=0.70\linewidth,height=0.52\textheight,keepaspectratio}}
% 图就地排版（float 宏包的 H）：图紧跟其说明文字，用户可边看图边照做。
% 房式原件设的是 tp（图只放页顶/独立图页）——那会让图浮到所属小节标题之前，
% 安装类文档里对照起来容易错行，故此处偏离。改回 tp 即可还原房式行为。
\makeatletter\AtBeginDocument{\def\fps@figure{H}}\makeatother

% 正文中文字体、中文断行规则等只对中文版生效的部分见 build/a3_cjk.tex，
% 由 build/build.sh 追加在本文件之后。

